\documentclass[11pt]{article}
\usepackage[margin=1in]{geometry}
\usepackage[T1]{fontenc}
\usepackage{lmodern}
\usepackage{amsmath}
\usepackage{booktabs}
\usepackage[colorlinks=true,urlcolor=blue,citecolor=blue]{hyperref}
\title{Testing Predictive Commitments in Interactive Reasoning\\\large Research draft for ARC Prize 2026}
\author{Alexander Serov\\Independent Researcher}
\date{October 29, 2026}
\begin{document}
\maketitle
\noindent\textbf{Status: research draft; development pilot failed feasibility.}
All four development arms have run. P failed the prespecified prediction-packet coverage threshold, leaving the matched efficacy comparison ineligible. At this revision's prelaunch cutoff, the main audit and calibration remain unexecuted. The full published-policy reproduction was attempted and failed. The corresponding P entry now has official public score 0.00, with rerun traces and resource use unavailable. The earlier bounded engineering entry is a separate result.

\begin{abstract}
An interactive reasoning agent can make useful predictions and still spend its inference budget poorly. We implemented a four-arm ARC-AGI-3 pilot to test whether recent action-effect prediction errors help allocate shared inference effort. The primary comparison holds the prediction interface and probability/outcome moments fixed, adding a pairing-sensitive Brier-error representation to a learned progress model under fixed regularization. A prospectively corrected collection supplied 16 eligible development windows; only three endpoints were positive, all from one game. B1 and A1 met the prespecified 50\% valid-packet coverage requirement. P supplied valid packets for only 89 of 417 eligible actions (21.3429\%), so the frozen pilot failed feasibility and provides no eligible efficacy comparison. The corresponding P entry received official public score 0.00; its actual rerun traces and resource use remain unavailable. A full published-policy reproduction attempt also failed, leaving both complete-run aggregates unavailable. The main audit and held-out progress calibration remain unexecuted at this revision's cutoff. This is a feasibility-failure pilot with an implemented public method, not evidence of allocator benefit or verification necessity.
\end{abstract}

\section{Question and scope}
We ask whether a fixed representation of observable prediction errors improves inference allocation beyond existing progress signals and matched confidence/outcome history. The intended contribution is a controlled test of that incremental signal. Predictive error need not imply that more computation is useful: it can mark a mistaken explanation, an unpredictable event, or an observation that the interface cannot label. A learned allocator must therefore be judged by subsequent progress and actual resource use, rather than by the mere presence of prediction checking.

ARC-AGI-3 permits an agent to act, inspect public observations, and revise its approach while sharing finite inference resources across games. Our study concerns that allocation problem. The paper prize values accuracy, universality, progress, theory, completeness, and novelty equally, and requires a working Kaggle entry\cite{arc2026}. The present manuscript reports an implemented method and a development feasibility failure; novelty and scientific effectiveness remain unverified.

\section{Baseline and closest mechanisms}
Our selected reproduction baseline is Daniel Franzen's public ARC-AGI-3 Milestone 2 notebook\cite{franzen}, built on Tufa Labs' Duck harness\cite{duck}. The full reproduction attempt failed; the published baseline score has not been reproduced. Its existing scheduling and action safeguards are essential comparators. Generic priority scheduling, memory, multi-agent critique, or a further action guard will not establish a new contribution. We preserve the upstream notebook configuration rather than substituting the repository's experimental defaults. Licensing and hardware requirements remain part of the reproduction audit.

Rodionov's plan executor compares predicted and observed settled frames after nonterminal actions and interrupts mismatching plans; direct game-client access can bypass it\cite{rodionov}. Kepler also uses retrospective transition checks and conditional prediction checks, which motivate cost-conditioned and verification-aware reporting\cite{kepler}. Active inference uses predictive hypotheses and information gain to select informative actions\cite{activeinference}. We make no novelty claim for these concepts, and squared error is not an information-gain estimator.

A public project plan already proposes shared inference priority conditioned on prediction accuracy and confidence calibration, with completion within a future window as its target\cite{dcwplan}. We inspected the committed plan, not its implementation or experimental validation. Given this close conceptual overlap, our question is limited to the controlled incremental value of the paired-error coordinate under the frozen representation and resource policy. A separate four-variant study reports differences that depend on model and inference effort; its bundled verification treatment consumes more cost-weighted tokens, while held-out performance remains untested\cite{rodionovablation}. These findings motivate matched controls and do not establish that verification is necessary.

\section{Proposed intervention}
\subsection{Commitments and observable errors}
For an observation history $h_t$ and a proposed action $a_t$, the interface accepts an optional probability vector $q_t$ before execution. A supplied vector forecasts three events of that next action: whether the frame changes, whether a level completes, and whether the attempt fails. Omission is a permitted abstention. A host evaluator binds the immutable commitment or abstention to the exact action and payload, records dispatch and return, and scores a valid supplied prediction once against the subsequent permitted public observation $o_{t+1}$. Later explanations cannot revise that commitment. The initial interface admits one predicted action per dispatch; an explanation-revision controller is outside this study.

The evaluator uses the action's last returned public frame. Public completed-level counts and failure metadata supply event labels where available. A refused action has no executed target. RESET and undo remain recorded transitions but are excluded from the initial prediction-error sample; an automatic reset cannot replace the preceding action's target. When the interface withholds a label, the evaluator records it as unavailable. Environment source, hidden state, future observations, and game-specific solution rules are excluded.

For event $e$, a valid probability $p_{t,e}$ and observable binary outcome $y_{t,e}$ give squared error $(p_{t,e}-y_{t,e})^2$. The history retains the last eight eligible executed actions and averages only the valid, observable pairs for each event. An empty pair set yields a missing residual. Coverage and outcome availability remain separate common features, so missing predictions cannot silently become correct predictions. The current allocator is event-only; a visual-mask feature has not been adopted for this study.

\subsection{Progress prediction and allocation}
The action-effect forecasts describe what the next action may do; the learned progress model addresses where to allocate inference effort. The learned arms use the same progress features: baseline priority, generated tokens, completed levels, action history, observation novelty, and prediction availability. Both also receive the mean probability, squared probability, and outcome for exactly the pairs used in each event's error. P adds the event Brier means; A1 omits them. Neither arm receives a game identifier, a solution table, or future observations.

The fitted progress probability $\hat r_i$ estimates whether game $i$ will complete its current level within a bounded future window. Each training window stores features at a scheduling boundary and labels completion before the first of 8,192 subsequent target-generated tokens or 20 subsequent real nonautomatic actions. Completion on the 20th action is positive unless the token horizon closed the window earlier. Reaching a horizon without completion is negative. Timeout, crash, or a resource stop leaves an unfinished window censored; it does not create a negative label. These observational labels describe progress under the collection policy, not the causal benefit of additional computation.

The development collection used uniform round-robin allocation over the four development games and retained all 16 planned slots. Its frozen rules require every opened window to be complete. Only a source-proven contiguous never-opened suffix after public WIN may be exempt, following an entirely complete opened prefix. Collection-level integrity, source, clock, resource, or cleanup failure blocks all fitting. Both endpoint classes are required, with identical rows and equal row weights in A1 and P.

The progress models are logistic classifiers with L2 coefficient 1.0, an unpenalized intercept, 2,000 full-gradient steps, and learning rate 0.05. Standardization and missing-value rules use development data only. A model is not updated on evaluation outcomes.

After initial service, allocation proceeds at admission boundaries. The scheduler snapshots the settled history, computes the progress probability, and freezes the features and learned multiplier for that ticket. A1 and P multiply the upstream baseline priority for game $i$ by $0.5+\hat r_i$. The common admission capacity is two complete request/tool exchanges. The upstream baseline priority may refresh while the ticket waits, so its final priority can change without recomputing $\hat r_i$. When all event residuals are missing, P uses the predeclared A1 fallback and reports its frequency.

\section{Evaluation design}
The primary contrast is P versus A1, using the same model, prediction interface, instrumentation, and common feature rules. Within each history, probability/outcome moments and Brier error use the same scored pairs. For binary outcomes,
\[
\overline{(p-y)^2}=\overline{p^2}+\overline{y}-2\overline{py}.
\]
Given the common moments, Brier and the paired interaction span the same linear predictor space, but fixed L2 shrinkage depends on the feature basis. Any benefit would concern this pairing-sensitive representation under fixed regularization; it would not establish a special advantage of the Brier coordinate or verification necessity.

The mandatory constant-probability controls have an exact representation-level interpretation. Replacing probabilities by $q$ on the same scored pairs gives $b_q=q^2+(1-2q)\bar y$. Our preprocessing centers observed values, maps missing values to zero after standardization, and shares all missingness indicators between arms. At $q=1/2$, each added Brier column is therefore zero. For $q\ne1/2$, it is a signed duplicate of the standardized paired-outcome column whenever that outcome varies, and zero otherwise. These controls add no pairing information; varying probabilities can still reveal $\overline{py}$ beyond the common marginals.

Duplication can nevertheless change regularization. For an isolated standardized predictor $z$ with effective coefficient $\beta=u+s v$, $s\in\{-1,1\}$, minimizing $\lambda(u^2+v^2)/2$ at fixed $\beta$ yields $\lambda\beta^2/4$, half the single-coordinate penalty. This calculation does not halve the penalty of an entire correlated design, and our fixed-step optimizer can also respond to reparameterization. The controls diagnose representation sensitivity rather than policy performance; they predict neither the magnitude of a fitted difference nor allocator benefit. They require no additional performance fleets.

\begin{center}
\small
\begin{tabular}{@{}p{0.08\linewidth}p{0.42\linewidth}p{0.42\linewidth}@{}}
\toprule
Arm & Inference allocation & Role \\
\midrule
B0 & Baseline prompts and batch policy under common resource accounting and admission & Practical baseline; distinct from published-policy reproduction \\
B1 & Shared prediction interface and instrumentation; baseline priority & Control for the total instrumentation treatment \\
A1 & Learned multiplier from common progress, coverage, and paired-moment features & Control for learned allocation and marginal confidence/outcomes \\
P & Same as A1, with event Brier means added & Primary fixed-representation test \\
\bottomrule
\end{tabular}
\end{center}

The contrasts answer different questions. P versus A1 tests the added pairing-sensitive representation, so P versus B1 cannot replace it. B1 versus an actual B0 can describe the total prompting, parser, single-action, and instrumentation treatment under the common resource policy; it cannot isolate prediction overhead. Actual matched B1 and B0 runs remain required for the later practical comparison; published-policy reproduction is separate.

The development fleet checks completeness and feasibility under the frozen resource policy; it does not replace the audit. The unexecuted main audit fixes ten public games, LLM request seeds 101, 202, and 303, all four arms, resource ceilings, and failure rules before its outcomes. Its prospective primary endpoint is observed public-prefix progress under the fixed limits, using the selected SDK's per-episode scalar formula. It includes fully observed \texttt{won}, \texttt{gave\_up}, \texttt{cancelled}, and \texttt{crashed} terminations without inferring why they ended.

For completed level $i$ with positive executed action count $a_i$ and public baseline $b_i$, define $s_i=\min(115,100(b_i/a_i)^2)$; otherwise $s_i=0$. With weights $w_i=i+1$ over all levels, the episode score is $\min(\sum_i w_i s_i/\sum_i w_i,100\sum_{i:s_i>0}w_i/\sum_i w_i)$. The fleet mean retains all ten games. Calculation requires a fully joined original public prefix and all source, cost, native, clock, cleanup, and original packet-feasibility conditions. Any unavailable game leaves the fleet mean unavailable. Raw framework states and stored scores remain separate; a missing stored scalar may be reconstructed only from complete joined observations. This prospective endpoint does not rescore the original DEV fleets or establish an authenticated SDK or official Kaggle result. The adopted methods are retained in \texttt{research/tier2-audit-protocol-20261028.md}.

A shared fleet scheduler creates interference across games, making game-level differences descriptive. Planned fresh fleet runs use matched LLM request seeds to characterize sensitivity on this fixed public sample. The entry does not pass these seeds to the SDK game constructor. Request seeds alone do not establish independent repetitions, deterministic GPU execution, or generalization to unseen games. Public games are exposed to previous research and potentially model training, so they provide development evidence rather than an uncontaminated hidden test.

The prespecified Tier 1 feasibility condition requires valid event predictions on at least 50\% of eligible executed actions separately in B1, A1, and P, with a nonempty denominator. Missing or malformed prediction packets abstain and remain in the denominator. RESET/undo and known refusals are excluded; missing or duplicate receipts and unknown execution invalidate completeness. Valid-packet coverage is distinct from event-label availability: report each event's available outcomes, scored pairs, and Brier denominator separately, retaining missing labels as unavailable. A valid packet with unavailable labels still counts toward packet coverage. This feasibility check does not establish accuracy superiority.

Resource ceilings do not guarantee equal actual computation. Report complete accepted target costs, all known retries, summaries, tails, failures and setup costs, with supported input/prefill and cache measurements and unavailable dimensions stated explicitly. Main-purpose output can include discarded replies and tails; it is charged in full without a causal token attribution. Unknown terminal target usage disqualifies the paired scientific comparison. Unsupported draft-work or GPU measurements remain unavailable, not zero. Scores and failures must be reported with the observed resource vector; equal quotas do not imply equal consumption or efficiency.

The future methods amendment replaces the unidentifiable isolated prediction-output ceiling. Generated Python can compute or reuse a prediction object, so its serialized size does not identify the incremental generation cost of eliciting it. The original 15\% gate remains unmeasured and its collection freeze unchanged. The future fixed-interface comparison makes no isolated prediction-overhead claim.

\section{Results: a failed feasibility pilot}
The four development fleets ran, but P failed the prespecified prediction-packet feasibility condition. The unchanged mechanism therefore has no eligible P-versus-A1 efficacy comparison. The other arms' evidence checks and the retained target-cost records cannot replace that condition. The official toolkit\cite{sdk} supplies the environment interface; a deterministic CPU smoke policy validates loading, permitted actions, trace capture, and artifact creation without providing a competitive baseline. Earlier engineering observations and their resource scopes are retained in Appendix~\ref{app:engineering}.

The corrected B0 development run completed on October 28 and passed its frozen evidence, accounting, and cleanup checks. Its 215 observed requests retained 446,084 accepted target-output tokens: 427,632 during gameplay and 18,452 during late drain, with no observed unknown or pending attempts. One game completed all six levels. The other three retained incomplete or crashed terminal records, so the descriptive four-game score is unavailable; the completed game's score cannot substitute for that aggregate. This is one practical-baseline fleet on exposed public games, not an official method score. Claim O22 retains the original readout and missing-score decision.

B1 passed the same evidence checks and supplied valid predictions on 331 of 422 eligible actions. Its 287 known requests retained 412,843 accepted target-output tokens: 408,789 during gameplay and 4,054 during late drain, with no observed unknown or pending requests. One game won all six levels; cancelled or crashed records leave the original four-game aggregate unavailable (O23).

A1 supplied valid packets on 276 of 374 eligible actions. Its first reader invocation refused a tuple-versus-JSON-list comparison of model parameters. A disclosed one-line format correction normalized the expected payload through JSON before the same equality check; one separate adjudication of the unchanged original evidence passed all 21 checks. The original refusal remains retained. No GPU replay, refit, model, runtime, budget, or endpoint change occurred (O24--O25). A1 retained 390,586 accepted target-output tokens across 330 known attempts: 388,225 during gameplay and 2,361 late, with zero observed unknown or pending attempts. All four raw framework states were \texttt{crashed}, and its original-rule four-game aggregate remains unavailable.

P supplied valid predictions on 89 of 417 eligible actions, or 21.3429\%, below the unchanged 50\% threshold. The remaining 328 actions had \texttt{prediction\_not\_supplied}; the complete diagnostic recorded zero invalid predictions and zero unknown transitions. All other 20 reader checks passed, but overall scientific eligibility remains false. P retained 404,393 accepted target-output tokens across 303 known requests: 398,009 during gameplay and 6,384 late, with zero observed unknown or pending requests in that scope. No efficacy score or surviving-subset comparison is reported (O26). This failure concerns prediction supply under the frozen optional interface; it does not estimate the effect of the added Brier coordinates on progress. The omissions do not establish model intent or unsupported keyword placement as their cause.

The corresponding unchanged P entry received official Kaggle public score \textbf{0.00}, submission ID 57749129 (O31). The submission was accepted at 21:28:14 UTC on October 28. Its original observer received API \texttt{COMPLETE} and the public score at 2026-10-28T21:42:51.801722 UTC; this is an observation time, not an exact server completion time. The request explicitly selected kernel 141404446, version 1. Final local custody matched the submission and requested version, but the server did not echo the version. No Submit retry occurred. Actual competition-rerun journals, runtime, token use, and private score remain unavailable. This documents a corresponding scored P entry, without establishing an action trace, allocation efficacy, novelty, or model/prize approval. It is distinct from the earlier B1/C1 engineering submission 57087987, also scored 0.00.

The full published-policy reproduction attempt (Bpub, kernel 141407495/version 1) retained the original 25 public games, four passes, 100 planned owners, and legacy C10 scheduling. Its observer recorded API \texttt{ERROR} at 22:10:23.061729 UTC on October 28; complete original retrieval finished at 22:12:53.254375 UTC, preserving 532 files and 6,907,564 bytes. The first completed strict readout was ineligible. Both full published-policy and SDK0.9.8 best-of-four-per-game aggregates remain null; the upstream score was not reproduced (O32). All 100 planned slots remain in the readout, without a subset average or score rescue.

The scan retained 100 host journals, four unknown-usage attempts, and zero observed pending attempts. Complete accepted target tokens, whole-kernel GPU use, and draft work remain unavailable: the known-token lower bound of zero is not a measurement of zero cost. The retained log records the 720-second benchmark release while the model was still loading. Four generic \texttt{TransportError} attempts with unknown usage preceded fail-closed accounting and \texttt{session\_binding\_failed} propagation. The underlying operating-system exception was not preserved; an unavailable listening service is consistent with the evidence, not a proven transport cause. The added accounting wrapper raised before upstream retry handling and before recording the original exception; this attempt therefore does not establish the behavior of the unmodified upstream retry policy. This is a failed reproduction attempt, separate from official P's score and the DEV packet-feasibility failure.

The failed collection and a separate official engineering entry have distinct evidential roles. The collection remains ineligible for fitting, and the engineering score does not evaluate the proposed allocator. Their complete counts, request associations, and accounting limits are retained in Appendix~\ref{app:engineering}.

A new collection used a prospectively frozen JSON transport correction while preserving the original scientific windows, endpoints, optimizer, and failure rules. Its complete 261-file output passed the frozen reader's source, clock, resource, cleanup, and original-slot checks. One direct replay-and-fit invocation retained all 16 slots and trained both models on 16 complete windows: 13 negative and 3 positive endpoints. All positives came from one game and episode. The collection-owned journal retained 153,147 known accepted target-output tokens in 112 usage records and 120 executed-action receipts, with no observed unknown usage or execution in that scope. Whole-kernel and draft-model work remain unavailable. Evidence claim O15 records the exact inputs and outputs; the failed predecessor remains ineligible.

The A1 and P design matrices both had rank 14, with 63 and 66 columns respectively. Half of the windows had no scored pairs for each event. The constant-half representation control produced exactly identical fitted training probabilities, while the constant-quarter and three-quarter controls each produced a maximum absolute difference of 0.0565. These are in-sample representation diagnostics under fixed regularization, not policy results. The small sample, concentrated positive labels, missing residuals, and absence of held-out calibration limit interpretation. The failed inference feasibility condition prevents a test of whether P improves on the matched A1 control.

The same collection provides little support for a mask extension. Only one recorded transition supplied a valid mask with an observable residual, $2/3$. It appears in the last-eight-action histories at two of the 16 window openings; the other 14 have no mask residual. With development-mean imputation, scale one for a constant column, and no extra P-only missingness indicator, the additional numeric mask column is zero on every training row. A zero-initialized coefficient therefore stays zero under the fixed L2 gradient updates, leaving the other coordinates unchanged mathematically. This is an analytic consequence of sparse development support, not a second executed fit or evidence that masks lose held-out validation. Claim O20 retains the exact opening joins and source counts. Five-game progress calibration remains unperformed.

\section{Reproducibility, disclosure, and limitations}
The evidence registry in \texttt{paper/claims.json} separates hypotheses, external reports, local observations, and official Kaggle results. Source revisions, licenses, model versions, configurations, raw-journal checksums, and failed predecessors accompany each observation. The draft, selected full method source, fixed allocation parameters, and prospective study documents are public at \url{https://github.com/AGoldian/predictive-commitments-arc2026}. Private notebook downloads and account metadata are excluded. The exact full method source was public before the official P request. Its Kaggle notebook remains private under the competition-source embargo until the competition ends. The reported score does not resolve all model and contest-eligibility obligations or supply the missing rerun evidence.

The pilot's conclusion is a feasibility failure of the frozen P policy. B1 and A1 exceeded the packet-coverage threshold, while P did not; this single sequential cohort does not identify why their coverage differed or establish an effect on allocation quality. The original DEV aggregates remain unavailable under their original endpoint. The prospective main audit and five-game calibration remain unexecuted at this revision's prelaunch cutoff. The full Bpub attempt failed its strict complete-run readout, with unknown target usage and no reproduced aggregate. Neither that failed reproduction nor the separate official P score changes the DEV feasibility refusal. Their frozen plans are not completed evidence, and the public-prefix endpoint cannot be applied retrospectively to repair DEV. No deployment benefit, efficacy advantage, or isolated prediction-overhead estimate follows from these observations.

Separate LLM advisors examine novelty, methods, theory, reproducibility, and scientific language. Each concern must reference evidence and retain uncertainty. Role diversity within one model does not imply independent model families. Known LLM judging biases limit the interpretation of reviewer scores\cite{judge}; averaging opinions cannot resolve an open empirical blocker.

LLM assistants are used for research discovery, code, critique, and drafting. Human authors must verify the sources, artifacts, and final text. Free and local resources are the current resource policy; full baseline inference needs compatible NVIDIA hardware. The final paper must record asset licenses, actual model versions, resource use, author details, and the code submission identifier. Public-development exposure, the small fixed sample, policy shift between collection and learned allocation, feature-basis dependence, and unavailable cost dimensions constrain any later interpretation.

This manuscript is licensed under Creative Commons Attribution 4.0 International (CC BY 4.0). Original research code and fitted allocation parameters use MIT-0; upstream code and model assets retain their separate terms.

\appendix
\section{Engineering observations and their scope}\label{app:engineering}
\subsection{Failed collection before the transport correction}
An October 28 four-game collection failed its frozen eligibility checks despite API \texttt{COMPLETE}. After a disclosed source-profile correction in the reader, its source-owned journal recorded 256 tool invocations: 85 clean sandbox completions, 3 spawned and 167 pre-spawn JSON-node resource failures, and 1 ordinary pre-sandbox return. All 16 planned slots were retained; the collection was ineligible for fitting. Collection-owned ledgers retained 148,912 known accepted target-output tokens across 260 usage records. They also retained 87 executed-action receipts, with zero observed unknown usage or execution in that journal scope. Kernel-wide token accounting and draft-model work were not measured. These observations identify a pipeline failure, not allocator performance. Evidence claim O14 retains both reader refusals and the complete failure audit.

\subsection{Separate official engineering entry}
A separate bounded private B1/C1 engineering entry received official Kaggle public score \textbf{0.00}, submission ID 57087987. An exact-reference API observation started at 00:11:22 UTC on October 12 and returned \texttt{COMPLETE}; the timestamp is not an exact server completion time. The request explicitly selected saved notebook version 1, while the server did not echo the version. The result and request association were independently checked against retained receipts. Raw competition-rerun journals, runtime, and usage remain unavailable. This score does not evaluate the proposed allocator; publication and prize eligibility of that engineering entry remain unresolved. Evidence claim O13 records these limits.

\subsection{First practical-baseline inference attempt}
The first matched-resource B0 attempt ended in Kaggle \texttt{ERROR}. Its complete 74-file output was retained, and one invocation of the frozen inference reader refused eligibility. The worker requested a stop and was killed; the Supervisor recorded complete physical cleanup, but no final worker-result document was written. A costs-first scan retained a provisional lower bound of 431,745 known accepted target-output tokens across 229 known attempts and two unknown terminal attempts. The reader also incorrectly required the pre-install parent environment to match the separate serving environment; the actual serving versions matched the prior successful collection. This attempt supplies neither a valid comparison nor a zero score. Claim O16 retains the original refusal; source-backed deadline diagnosis and any future correction are separate.

Three subsequent B0 corrections also failed before an eligible comparison. One exceeded the 25-minute model-startup limit. After a prospectively increased startup allowance, the next printed READY after 1,552 seconds but failed at its first physical-request hook, retaining one unknown-usage attempt. The following package failed before model setup because its embedded runtime/profile hashes were stale. Claims O17--O19 preserve each original refusal, diagnosis, and stopped sequence; no remaining arm was launched by those sequences. Missing host streams do not establish zero whole-kernel or GPU cost. The current packaging correction selects the exact packets accepted by a single local boot-control check; that check establishes source association, not allocator performance.

\subsection{Baseline serving and gameplay}
A bounded GPU infrastructure pilot served the pinned target/draft model and recorded 14 actions including RESET on one public game in 120 seconds, completing zero of seven levels. Priority scheduling was disabled. This verifies model serving and gameplay in that configuration, not full baseline reproduction. Action-attributed, analyzer, and server token counters disagree; their scopes must be reconciled before resource comparisons.

A second bounded baseline run on three selected development games completed 6/6 levels on ft09, 2/7 on ls20, and 3/7 on vc33, with 517 framework-history actions. The latter two games were cancelled at the gameplay ceiling. Completed API responses reported 486,120 output tokens, including 4,001 after ls20's last recorded action. These observations establish longer baseline gameplay on this fleet; they do not evaluate the proposed intervention. A source-level audit found that three games used at most three of ten gate slots and that analyzer yield did not release a slot. At that stage, a resource-normalized admission boundary still had to be implemented and reviewed before allocation comparisons.

\subsection{Admission capacity and lifecycle}
A subsequent three-worker capacity diagnostic retained 38,488 known accepted target-output tokens: 37,906 observed within its fixed 300-second interval and 582 during late drain. No physical attempt had unknown usage in the retrieved ledger. The diagnostic failed its frozen admission criterion because the first automatic warmup RESET occurred outside a complete logical lease. Its eligible throughput is therefore unavailable; known costs remain charged. Notebook completion did not establish a feasible capacity or authorize a scientific comparison.

A distinct prospective RESET-prefix correction subsequently passed its frozen analyzer, retaining 32,991 known tokens across 80 physical requests: 31,941 within the fixed interval and 1,050 during drain, with no unknown usage. Independent replay reproduced the identical analysis. This one-seed engineering observation does not evaluate the proposed mechanism or establish an optimal capacity.

The subsequent capacity-two RESET-prefix run also passed its frozen analyzer and independent replay. It retained 65,041 known accepted target-output tokens across 105 physical requests: 64,454 observed within the fixed 300-second interval and 587 during late drain, with zero unknown-usage attempts. The capacity-one and capacity-two observation-time proxies were 106.47 and 214.8466667 accepted tokens per fixed second, respectively. Their ratio, 0.4955627, fell below the frozen 0.8 threshold. This selected capacity two for prospective common admission between complete request/tool exchanges without using game outcomes. These were one-seed development engineering runs executed in capacity-one then capacity-two order; the proxy is not generation speed, accuracy, efficacy, or optimality, and no official score is established. All four matched arms use the chosen capacity in the prospective execution freeze. A complete matched comparison remains unavailable.

Separately, three owned CPU lifecycle cases passed their original strict replay checks for batch retention, cutoff-prefix drain, and external stopping. They used HTTP/SDK doubles, totaling 35 synthetic accepted tokens and five synthetic sends, alongside seven actual native tool subrecords. The cutoff retained its entered action and refused the remaining batch suffix; the external-stop case retained its expected worker termination. These observations support only the tested lifecycle composition. Synthetic tokens are excluded from actual model usage. The checks do not establish real-SDK/model compatibility, isolation against arbitrary hostile code, or allocation effectiveness. Configurations, failed predecessors, source hashes, and replay records are preserved in the accompanying evidence registry.

\subsection{Runtime isolation limits}
The inherited Python sandbox exposes a reflection route to unrestricted modules. Restricted imports alone do not establish isolation from environment files or unlogged network access. The pilot did not establish any use of this route. A separately patched host/bootstrap passed benign Kaggle CPU checks for arithmetic, retained-function restoration, and owned action receipts under mandatory native policy setup and declared process/resource bounds. These fixture checks do not certify isolation. Full real-model integration, cross-arm compatibility, and the final competition gateway lifecycle remain open.

\begin{thebibliography}{99}
\bibitem{arc2026} ARC Prize Foundation. \emph{ARC Prize 2026 Paper Prize}. \url{https://arcprize.org/competitions/2026/paper}. Accessed October 11, 2026.
\bibitem{franzen} Daniel Franzen. \emph{ARC-AGI-3 Milestone 2 solution}. \url{https://github.com/da-fr/arc-agi-3-solution}. Revision recorded in the artifact registry.
\bibitem{duck} Tufa Labs. \emph{Duck harness}. \url{https://github.com/Tufalabs/duck-harness}. Upstream source and Kaggle notebook.
\bibitem{sdk} ARC Prize Foundation. \emph{Official ARC-AGI toolkit}. \url{https://github.com/arcprize/ARC-AGI}.
\bibitem{judge} L. Zheng et al. \emph{Judging LLM-as-a-Judge with MT-Bench and Chatbot Arena}. 2023. \url{https://arxiv.org/abs/2306.05685}.
\bibitem{rodionov} Sergey Rodionov. \emph{Executable World Models for ARC-AGI-3 in the Era of Coding Agents}. 2026. \url{https://arxiv.org/abs/2605.05138v2}.
\bibitem{activeinference} Karl Friston et al. \emph{Active inference and artificial reasoning}. Nature Communications 17, 10125, 2026. \url{https://www.nature.com/articles/s41467-026-77209-5}.
\bibitem{kepler} Wensen Wu. \emph{Kepler: Auditable World Models for ARC-AGI-3}. 2026. \url{https://arxiv.org/abs/2610.00834v1}.
\bibitem{dcwplan} dcw06. \emph{ARC-AGI-3 Project Plan 7}, Section 8.16. Public design document, revision 54e129671e1583e958dbde2c34aede1a7d341b72, October 3, 2026. \url{https://github.com/dcw06/ARC-AGI-3/blob/54e129671e1583e958dbde2c34aede1a7d341b72/docs/ARC-AGI-3_Project_Plan_7.md}. Accessed October 12, 2026.
\bibitem{rodionovablation} Sergey Rodionov. \emph{Do Coding Agents Need Executable World Models, Simplification, and Verification to Solve ARC-AGI-3?} Version 2, August 27, 2026. \url{https://arxiv.org/abs/2607.15439v2}. Methods, design, and resource reporting inspected October 28, 2026.
\end{thebibliography}
\end{document}
